\chapter{连线世界}

\begin{branch}
\goodend 两人走出温泉旅馆，身体如褪去云层后的澄空般清澈。月亮已经从地平线上升起，宣告着美好的一天的终结。“已经这么晚了啊，”铃还回忆着在温泉池旁吃的那盘美味的炒面，“现在再坐地铁回去，到家未免会太晚了。我们坐出租车吧！”彩名接受了这个提议。于是在和司机大叔深刻地交换了对于俄乌局势的看法后，两人回到各自的出租屋中，结束了快乐的周末约会。铃想向艾莉丝分享今天神奇的中奖经历；可惜她还没有回来。
\end{branch}

再怎么说，坐出租车也太破费了。铃放弃了这个小布尔乔亚方案，和彩名走向了来时的地铁站。残存的日光和新生的月光混合着，从两个方向映照在即便无人也还在勤劳运行的扶梯上，不留下阴影。或许古希腊人在创作西西弗斯神话的时候就已经预料到了自动扶梯的发明。铃和彩名并排站在扶梯上；出于扶梯寿命考虑，市地铁现在已经不再提倡左行右立。她们离开无温度的日光和月光，而后被人造的同样无温度的荧光笼罩。

即使周末有晚高峰，现在也已经过了那个时候。站厅里只有寥寥几个路人，沿着他们预定的道路行进。安检人员显然在无聊地等待即将到来的预定的交接班时刻。他似乎想要履行自己的职责、仔细检查彩名那复杂的服饰，但他最终放弃了，只是拿金属探测棒随意地扫了扫，全程不发一言。直到两人快要走远，他才提醒似的，在她们背后自言自语，“地铁运行计划今日没有调整。”他接着敷衍地念出宣传词，“海汶地铁，连接世界。”

铃刷卡通过闸机。屏幕上显示卡片余额仅有个位数。“应该不会有问题，”铃通过了闸机，后者在她身后决然地关上，“大不了出站的时候再充值。”彩名牵起铃的手，“莉莉姆可以帮你刷哦……”她立刻想起了地铁的计费方式，“好像不行。”

再次劳烦西西弗斯，她们走下了站台。相比常见的岛式，这个偏远小站使用的是侧式站台，一面是轨道，一面是墙壁。铃看到一位施工人员正拎着一个油漆桶，踩着立在绿色防护垫上的脚手架，往墙上的塑料布泼洒亮黄色的颜料。塑料布似乎完全没有贴整齐，黄色胶带随意地粘贴在其上。它缠绕着自身，混乱地遮蔽着其后的巨幅画作。铃看不清那是一幅什么样的画；她只能看到一片被半透明塑料布覆盖的深邃黑色。施工人员察觉到有人走近，连忙向她们招手示意，仿佛她们的到来会扰乱施工的进程。“他又不是在钓鱼；施工也怕鱼被吓跑吗？”铃和彩名便走开了，在站台的另一侧等待着地铁。不过，为什么施工人员要戴着镶钻面具呢？

等待了许久，地铁还没有到来。铃环顾站台。长长的站台望不见尽头，被荧光灯均匀地照亮，却看不到其他人，只留她俩和工人在这寂静中。如果站台的某些区域没有旅客，那为什么还要亮灯呢？铃思考着这一优化能带来的能源节省。只听几声清脆的咔嚓声，站台其余区域的灯光开始闪烁，随后逐渐变暗。看来地铁公司确实想到了这项环保措施。铃为自己的创新性感到担忧。孤独的光线洒在她们身上，浮起一座黑暗中的孤岛。

不过铃还想到了另一种可能：她们错过了末班车，所以灯光才开始逐渐熄灭。她不认为这种可能性有多大，因为她们从进站到来到站台并没有过去多久。只是以防万一，铃扭头看了看站台墙壁上的时钟；时钟早已停了，表盘扭曲着，表针指向一个不存在的时间。“这大概是《记忆的永恒》的仿品。”铃对时钟放弃功利效用感到不满。她掏出手机确认了时间，以及还会有地铁班次到来的推断。

\begin{branch}
\badend 对于时钟停摆的根据，铃还有另一种推测。她的背后被重重地推了一下，正如预料。铃的小身板在巨大的力量下无法自持，只能无助地向前倒去、扑向无光的地铁轨道。最后，她的耳边只剩下地铁列车划过轨道的悲鸣、彩名的尖叫、以及嫌疑人X的狂言：“我们不需要新的神明！我们受够了！”

\goodend 然而全国各地的新开行地铁都安装了屏蔽门。铃只是狠狠地撞在了屏蔽门上，让防爆玻璃发出一声低沉的抱怨。嫌犯很快就被警察带走了。
\end{branch}

“车来了哦！”彩名提醒正在妄想的铃，拉着她走近了地铁车厢。晚班的地铁车厢十分冷清，不过也有零零散散的乘客；这是好事，意味着她们可以随意地挑选座位。随着大腿感受到冰凉的金属座椅的触感，她们戴上耳机，隔绝了VVVF的歌声。座椅微微地振动起来，车窗外的昏暗的站台飞快地离去，随之离去的还有施工人员工作的那面墙壁。窗外只剩下闪动的黑暗。

显示屏上闪烁着线路和站点信息，提供一种目的地的保证。“海汶地铁，连接世界！”彩名把屏幕上滚动的宣传词念了出来；看起来她很喜欢这句宣传词，“铃，你不觉得地铁不仅连接着地上的世界，它本身就是一个连线世界吗？”如果她指的是地铁使各站相互连接，这样的比喻并非没有道理。“是啊，要是没有地铁，我们回去还得坐好久的公交车。”铃应和着，稍微调低了耳机的音量，以便听清彩名接下来的话语。

然而，伴随着VVVF舒缓的乐声传进耳朵的，不是彩名的话语，而是车厢中其他乘客的窃窃私语：

“听说了吗……你听说了吗……”

“XX学院有人跳楼了……现场一片惨状啊……”

“好想去看看……空气力学的最新研究……”

“死了！都已经死了！我们还要期待什么……”

“神明会有继任者吗……”

“你收到消息了吗……你是选举人吗……你是候选人吗……”

“不会了……这个世界已经结束了……”

铃想要寻找这些碎语的来源。她努力地向车厢前后望去。无名的乘客们或者坐着、或者站着、或者依靠在扶手上。他们盯着手中的手机屏幕，他们的脸被手机屏幕的光芒照亮，铃却仍然看不清他们的面容。他们的手指飞快地按压着手机，仿佛那是他们与这个世界唯一的接口。手机发出规律的键盘打字声，搅动着车厢内凝固的空气：“听说了吗……听说了吗……就在这里……”

车厢的灯光闪烁了一下；这或许是列车驶过了接触网交接的中性绝缘区段。当灯光重新亮起，铃只看到一张张看向自己的脸庞。从下方溅射光线的手机屏幕照亮了他们的面容，却仅仅是模糊的混沌。脸庞上依稀可见的嘴角以一种不可思议的程度扭曲着，向上升起，露出其下漆黑的深渊。“听说了吗？！”他们喊叫着，嘴唇却不动丝毫。

铃怀疑他们是在询问自己，但是她不敢作答，只是微微颤抖着双手。这仅仅是一条普通的地铁；1X号东部快线，全长六十余公里，连接四大市辖区；有点太长了，不过还在正常的范围内。这仅仅是一组普通的列车；CCD5555动车组，8A编组，6动2拖，轴重16吨，簧下质量1.8吨；很好，没有误用小编组；接触网供电，直流1500伏，电机功率二百千瓦，最高运营时速120千米每小时；速度有点慢了，跑完全程可能要很久，不过应该也没人坐T2T；这里的T指的是terminal而非telomere。很好，自己还很正常；铃默念着这些事实；这个世界也还很正常。没有，自己没疯；铃默念着这些事实；这个世界也没疯。那么现在是谁在看着我？那么现在是谁在对我说话？这一定有个解释。一定是自己泡温泉泡出幻觉了；一定是这样。自己本来就身体虚弱；在高温度和高湿度的条件下，外周血管扩张、血压下降，电解质失衡、血容量减少，迷走神经反射过度激活；这些因素都有可能影响自己的思考。这种情况是正常的，回家睡一觉就好。可是现在眼前的究竟是什么？可是现在耳中的究竟是为何？

恍惚间，铃的肩膀传来熟悉的触感；彩名紧靠过来。她用力握住铃放在大腿上的手，停止了它的颤抖。彩名笑着望向铃，“我们来唱歌吧！”她在这个荒诞的场景下提出荒谬的主意，不知是否也符合铃对“正常”的定义，“来唱‘莉莉姆之歌’！”彩名的十指扣住铃的十指。她高声歌唱，唱着铃熟悉却从未听过的歌：
\begin{song}
世界最后的魔法少女\\
就是本小姐！\\
凌晨三点 大眼睛瞪圆\\
床单怪兽 在对我眨眼\\
它们说 月亮卡了壳\\
需要我去 用彩虹糖解锁！\par
\halfline%
变出魔杖 飞上天空\\
任务是给 星星挠痒痒！\\
粉色能量 呼啦啦发射！\\
把忧郁蘑菇 变成彩虹糖吃掉！\par
\halfline%
天使也好 恶魔也罢\\
全部都是 我的手下败将！\\
旋转 跳跃 我闭着眼！\\
坏蛋在哪 大概在梦里面？\\
管他呢 快乐就起舞！\par
\halfline%
世界最后的魔法少女\\
就是本小姐！\\
变出魔杖 飞上天空\\
任务是给 星星挠痒痒！\\
天使也好 恶魔也罢\\
反正我看不见！\par
\halfline%
给世界一个大拥抱\\
我要让 世界充满笑容\\
天使也好 恶魔也罢\\
一起来跳舞！\\
任务是给 星星挠痒痒！\par
\halfline%
嗯哼 世界又和平了\\
奖励自己 一瓶亮片汽水！\\
最后的魔法少女 要去睡觉觉…\\
晚安 我的多彩宇宙！
\end{song}


荒唐的歌词，铃忍不住这样想；但她却也跟着唱了起来，大概是被彩名热烈的歌声打动了吧，抑或是因为彩名可爱的笑容。怪异的乘客仍然盯着她们，仅仅是盯着她们，不知道有没有听到两人的歌声。当整首歌唱完，彩名像是想起了什么，“呀，该吃彩虹糖了！”她掏出随身携带的药盒。铃感到羞愧；提醒彩名吃药本是她的职责。

“战斗的时刻到了！”彩名吃下彩虹糖之后，宣告了战斗的开始。铃也跟随她站了起来。彩名举起手中的长伞。铃担心她这样做会戳到车厢里的其他人；不过这样的忧虑是多余的，车厢内的“乘客”看起来可不像正常人。他们被彩名的长伞吸引了目光（如果他们有眼睛的话），停止了任何话语和动作，仿佛陷入了恐惧。列车仍在前行，转向架传来哐哐的震动，彩名手中的长伞却像被钉在了空气中。接着，彩名中气十足地念出一串魔法咒语，声音穿透VVVF的杂音：“立刻绝育所有流浪猫！”

铃不明白为何这样的魔法咒语能够生效；但它确实生效了。她也不明白这样的咒语包含着怎样的魔法；但它确实包含着魔法。在咒语被说出口的瞬间，车厢内的气氛瞬间转变了，一波波汹涌的巨浪推过空气。一半的“乘客”的脸庞瞬间展露狂喜，另一半的“乘客”的面容刹那间传为愤怒。他们似乎被彩名的魔法咒语迷住了心窍，不再关注这两位连线世界的闯入者，反而开始互相攻讦：

“她在说什么？！”

“你其实赞同这样的观点吧！我早就看出来了！”

“猫猫这么可爱，怎么能伤害猫猫！”

“流浪猫吃了多少鸟儿！城市的生态都被破坏了！”

“还不是你们没良心的人弃养导致的！”

“现代人怎么能离开猫呢？没有猫摸，人会死的！”

“我们都被猫奴役了！你知道吗，猫其实是从外星来统治地球的！”

彩名按下伞柄上的按钮。黑色的伞面倏地张开，伞骨划破空气，发出清脆的声响。强大的力量击倒了正前方的敌人。铃第一次看见伞面内侧的复杂的赤红纹路，像是魔法阵，也像是印刷电路板，隐隐地散发光芒。“一起来！”彩名的另一只手牵住铃，拉着她沿车厢向前进击。锐利的伞尖轻易地刺穿一个个扭曲的身影；近旁的身影也恐惧着，坍缩着消失在空气中；更远的身影却还是在讨论是否应该绝育流浪猫。铃这才发现这些“乘客”并非实体，而只是言语聚成的影像。彩名欢快地旋转、跳跃，挥动长伞扫过周围的敌人，只留下一阵阵尖叫作为伴奏。铃被她牵引着，行走着舞步。两人在天旋地转的车厢中跳起致命的双人舞。这绝非乘坐地铁的正确方式，铃心想；不过现在正确与否已经无所谓了；铃只是荒唐地感到开心。“保持笑容！”舞蹈教练莉莉姆给出指导。她们就这样笑着、舞着，在“乘客”的伴奏下走过一节节车厢。当她们最终到达最后一节车厢（细心的读者应该记得总共有8节车厢），彩名收起了伞，扣住铃的双手，舞完了最后一圈。

\begin{figure}[p]
\centering\noindent%
\makebox[0pt][c]{%
\includegraphics[width=0.8\paperwidth]{supp4.jpg}%
}%
\caption{{\heiti {\large 遮蔽·施工中}\hfill 创作者：画匠}\\
本装置作品以透明塑料布、黄色胶带、脚手架与施工工具构建出一个“尚未完成”的现场。通过刻意暴露的遮挡与临时结构，作品探讨了艺术创作背后的无名劳动、保护与脆弱、秩序与混乱的张力，以及人类在过渡状态中的不安与希望。它既是一个暂停的施工现场，也是一面关于信任、屏障与修复的隐喻之镜。}%
\end{figure}
